\documentclass[12pt]{article}
\usepackage[utf8]{inputenc}
\usepackage[T1]{fontenc}
\usepackage[margin=1in]{geometry}
\usepackage{setspace}
\usepackage{xcolor}
\usepackage[colorlinks=true, urlcolor=blue!60!black, linkcolor=blue!60!black]{hyperref}
\setlength{\parindent}{2em}
\setlength{\parskip}{1.25em}
\renewcommand{\baselinestretch}{1.0}
\begin{document}
\begin{center}
{\LARGE \textit{Egbert v. Boule}, 596 U.S. 482 (2022)\par}
\vspace{0.8em}
{\large Case Brief --- Four-Page Edition\par}
\end{center}
\vspace{1em}

\section{Memory Jogger}

An innkeeper roughed up by a Border Patrol agent in his own driveway could not sue under \textit{Bivens}: the Court rejected his Fourth Amendment claim 6--3 and his First Amendment retaliation claim 9--0. A court may not extend \textit{Bivens} if there is any reason to think Congress is better equipped to create a remedy (\href{https://supreme.justia.com/cases/federal/us/596/21-147/}{Justia}).

\section{Detailed Case Facts}

Robert Boule ran the Smuggler's Inn, a bed-and-breakfast abutting the Canadian border in Blaine, Washington, and also served as a paid informant for Border Patrol Agent Erik Egbert. In March 2014, after Boule told Egbert that a Turkish guest was arriving, Egbert followed the guest's car into the inn's driveway. When Boule asked him to leave, Egbert allegedly lifted him off the ground and threw him against the car and then to the ground, injuring him. After Boule complained to Egbert's supervisors and filed a Federal Tort Claims Act claim, Egbert allegedly retaliated by reporting Boule's ``SMUGLER'' license plate to the state and prompting an IRS audit, which cleared Boule but cost him over \$5,000.

\section{Procedural History}

Boule sued Egbert personally under \textit{Bivens v. Six Unknown Fed. Narcotics Agents}, 403 U.S. 388 (1971), for Fourth Amendment excessive force and First Amendment retaliation. The district court granted summary judgment for Egbert; the Ninth Circuit reversed, allowing both claims, and denied rehearing en banc over the dissent of twelve judges. The Supreme Court granted certiorari but declined to consider whether to overrule \textit{Bivens}.

\section{Judicial Votes}

Reversed: \textbf{6--3} on the Fourth Amendment claim, \textbf{9--0} on the First Amendment claim (decided June 8, 2022).
\begin{itemize}
\item \textbf{Majority:} Justice Thomas, joined by Chief Justice Roberts and Justices Alito, Kavanaugh, and Barrett.
\item \textbf{Concurring in the judgment:} Justice Gorsuch.
\item \textbf{Concurring in part and dissenting in part:} Justice Sotomayor, joined by Justices Breyer and Kagan (concurring as to the First Amendment claim, dissenting as to the Fourth).
\end{itemize}

\section{Holding}

\textit{Bivens} supplies no damages remedy for either Boule's excessive-force claim or his retaliation claim. Collapsing the two-step \textit{Bivens} framework into a single separation-of-powers inquiry, the Court held that a federal court may not recognize a \textit{Bivens} action whenever there is \textit{any} reason to think Congress might be better equipped to create a damages remedy.

\section{Analysis of Opinions}

\textbf{Majority (Thomas).} Creating a damages action is a legislative function, and the Court had refused to imply a new constitutional tort eleven straight times since \textit{Bivens}, \textit{Davis v. Passman}, and \textit{Carlson v. Green}. The two steps of \textit{Ziglar v. Abbasi}, Thomas wrote, ``often resolve to a single question.'' For the Fourth Amendment claim, two independent reasons foreclosed relief. First, border security: the question is whether courts should authorize suits against Border Patrol agents generally, which implicates national security. Second, alternative remedies: CBP's internal grievance process forecloses \textit{Bivens} even though it offers no damages, no participation, and no judicial review. For retaliation, a plaintiff can allege a retaliatory motive for almost any adverse action, so such suits would chill legitimate conduct across the federal workforce. The majority declined to overrule \textit{Bivens} itself.

\textbf{Gorsuch (concurring in the judgment).} Weighing the costs and benefits of new causes of action is always Congress's job; the Court should forthrightly overrule \textit{Bivens}.

\textbf{Sotomayor (dissenting in part).} The excessive-force claim was ``materially indistinguishable'' from \textit{Bivens}. The Court rewrote \textit{Ziglar}, stretched national-security deference to agents operating up to 100 miles from any border, and treated a toothless grievance process as an adequate remedy.

\textbf{Subsequent development.} In \textit{Goldey v. Fields}, 606 U.S. 942 (2025) (per curiam), the Court unanimously applied \textit{Egbert} to deny a \textit{Bivens} remedy for a federal prisoner's Eighth Amendment excessive-force claim.

\section{Examples: Future Applications}

\begin{enumerate}
\item \textbf{Same-Side (No Remedy): An FBI Raid Far From the Border.} FBI agents break a homeowner's arm in a warrantless Houston raid. The facts look like \textit{Bivens}, but the agents will argue a new category of defendants and an inspector-general process---likely defeating the claim unless the court finds no meaningful difference from \textit{Bivens}.
\item \textbf{Same-Side (No Remedy): Retaliation by a Federal Inspector.} A restaurant owner who criticized a food-safety agency is hit with baseless citations. \textit{Egbert} forecloses a retaliation \textit{Bivens} claim; the remedy is for Congress.
\item \textbf{Opposite-Side (Statutory Cause of Action): A State Trooper.} The same force by a state trooper is actionable under 42 U.S.C. \S~1983, because Congress expressly created that cause of action.
\item \textbf{Opposite-Side (Prospective Relief): Enjoining a Federal Practice.} Residents can sue an agency head for an injunction against a written policy of warrantless backyard entries; \textit{Egbert} limits only implied damages against individual officers.
\item \textbf{Fence-Sitter: A Near-Replica of \textit{Bivens}.} Federal narcotics agents on a joint task force arrest and strip-search a tenant without a warrant. Whether the task-force assignment makes a meaningful difference from \textit{Bivens} turns on the facts.
\end{enumerate}

\section{Critique}

\textbf{Remedial asymmetry.} State officers face damages liability under \S~1983 for the same violations that federal officers now largely escape. \textbf{Executive self-immunization.} If any internal complaint process counts as an alternative remedy, agencies can shield their officers by maintaining one. \textbf{Candor.} As Justice Gorsuch observed, keeping \textit{Bivens} while making its test nearly impossible to satisfy invites costly litigation over hair-splitting ``new context'' distinctions.

\section{Key Quotations}

\begin{quote}
``Because our cases have made clear that, in all but the most unusual circumstances, prescribing a cause of action is a job for Congress, not the courts, we reverse.''
\end{quote}
\noindent --- Justice Thomas, for the Court, 596 U.S. at 486.

\begin{quote}
``While our cases describe two steps, those steps often resolve to a single question: whether there is any reason to think that Congress might be better equipped to create a damages remedy.''
\end{quote}
\noindent --- Justice Thomas, for the Court, 596 U.S. at 492.

\begin{quote}
``The Court goes to extraordinary lengths to avoid this result: It rewrites a legal standard it established just five years ago, stretches national-security concerns beyond recognition, and discerns an alternative remedial structure where none exists.''
\end{quote}
\noindent --- Justice Sotomayor, concurring in the judgment in part and dissenting in part, 596 U.S. at 505.

\end{document}
